% from ega1-10-fr.tex
\oldpage[I]{194}

espace sous-jacent (\hyperref[I.10.6.1-fr]{10.6.1}). Comme
$S_\lambda\to S_\mu$ est un monomorphisme de préschémas, on voit donc
d'abord que les produits
$X_\lambda\times_{S_\lambda}Y_\lambda$ et
$X_\lambda\times_{S_\mu}Y_\lambda$ sont identiques
(\hyperref[I.3.2.4-fr]{3.2.4}), puis que
$X_\lambda\times_{S_\mu}Y_\lambda$ s'identifie à un sous-préschéma fermé
de $X_\mu\times_{S_\mu}Y_\mu$ ayant même espace sous-jacent
(\hyperref[I.4.3.1-fr]{4.3.1}). Cela étant, le produit
$\mathfrak{X}\times_{\mathfrak{S}}\mathfrak{Y}$ est la
\emph{limite inductive} des préschémas usuels
$X_\lambda\times_{S_\lambda}Y_\lambda$~: en effet, on voit comme dans
(\hyperref[I.10.6.2-fr]{10.6.2}) qu'on peut se ramener au cas où
$\mathfrak{S}$, $\mathfrak{X}$ et $\mathfrak{Y}$ sont des schémas formels
affines. Compte tenu de (\hyperref[I.10.5.6-fr]{10.5.6, (ii)}) et de
l'hypothèse sur les systèmes fondamentaux d'Idéaux de définition de
$\mathfrak{S}$, $\mathfrak{X}$ et $\mathfrak{Y}$, on voit aussitôt que
notre assertion résulte de la définition du produit tensoriel complété de
deux algèbres (\hyperref[0.7.7.1-fr]{0, 7.7.1}).

En outre, soit $\mathfrak{Z}$ un $\mathfrak{S}$-préschéma formel,
$(\mathscr{M}_\lambda)$ un système fondamental d'Idéaux de définition de
$\mathfrak{Z}$ ayant $I$ comme ensemble d'indices,
$u:\mathfrak{Z}\to\mathfrak{X}$,
$v:\mathfrak{Z}\to\mathfrak{Y}$ deux $\mathfrak{S}$-morphismes tels que
$u^*(\mathscr{K}_\lambda)\mathscr{O}_{\mathfrak{Z}}
\subset\mathscr{M}_\lambda$ et
$v^*(\mathscr{L}_\lambda)\mathscr{O}_{\mathfrak{Z}}
\subset\mathscr{M}_\lambda$. Si l'on pose
$Z_\lambda=(\mathfrak{Z},
\mathscr{O}_{\mathfrak{Z}}/\mathscr{M}_\lambda)$, et si
$u_\lambda:Z_\lambda\to X_\lambda$ et
$v_\lambda:Z_\lambda\to Y_\lambda$ sont les $S_\lambda$-morphismes
correspondant à $u$ et $v$ (\hyperref[I.10.5.6-fr]{10.5.6}), on vérifie
aussitôt que $(u,v)_{\mathfrak{S}}$ est la limite inductive des
$S_\lambda$-morphismes $(u_\lambda,v_\lambda)_{S_\lambda}$.

Les considérations de ce numéro s'appliquent en particulier lorsque
$\mathfrak{S}$, $\mathfrak{X}$ et $\mathfrak{Y}$ sont localement
noethériens, en prenant comme systèmes fondamentaux d'Idéaux de définition
les systèmes formés des puissances d'un Idéal de définition
(\hyperref[I.10.5.1-fr]{10.5.1}). Mais on notera que
$\mathfrak{X}\times_{\mathfrak{S}}\mathfrak{Y}$ n'est pas nécessairement
localement noethérien (voir toutefois
(\hyperref[I.10.13.5-fr]{10.13.5})).
\end{env}

\subsection{Complété formel d'un préschéma le long d'une partie fermée.}
\label{subsection:I.10.8-fr}

\begin{env}[10.8.1]
\label{I.10.8.1-fr}
Soient $X$ un préschéma (usuel) localement noethérien, $X'$ une partie
fermée de l'espace sous-jacent à $X$~; désignons par $\Phi$ l'ensemble des
faisceaux cohérents d'idéaux $\mathscr{J}$ dans $\mathscr{O}_X$, tels que
le support de $\mathscr{O}_X/\mathscr{J}$ soit $X'$. L'ensemble $\Phi$
n'est pas vide (\hyperref[I.5.2.1-fr]{5.2.1},
\hyperref[I.4.1.4-fr]{4.1.4} et
\hyperref[I.6.1.1-fr]{6.1.1})~; nous l'ordonnerons par la relation
$\supset$.
\end{env}

\begin{lemma}[10.8.2]
\label{I.10.8.2-fr}
L'ensemble ordonné $\Phi$ est filtrant~; si $X$ est noethérien, pour tout
$\mathscr{J}_0\in\Phi$, l'ensemble des puissances
$\mathscr{J}_0^n$ ($n>0$) est cofinal à $\Phi$.
\end{lemma}

\begin{proof}
En effet, si $\mathscr{J}_1$ et $\mathscr{J}_2$ appartiennent à $\Phi$, et
si on pose $\mathscr{J}=\mathscr{J}_1\cap\mathscr{J}_2$,
$\mathscr{J}$ est cohérent puisque $\mathscr{O}_X$ est cohérent
(\hyperref[I.6.1.1-fr]{6.1.1} et
\hyperref[0.5.3.4-fr]{0, 5.3.4}), et l'on a
$\mathscr{J}_x=(\mathscr{J}_1)_x\cap(\mathscr{J}_2)_x$ pour tout
$x\in X$, donc $\mathscr{J}_x=\mathscr{O}_x$ pour $x\notin X'$ et
$\mathscr{J}_x\neq\mathscr{O}_x$ pour $x\in X'$, ce qui prouve que
$\mathscr{J}\in\Phi$. D'autre part, si $X$ est noethérien, et si
$\mathscr{J}_0$ et $\mathscr{J}$ appartiennent à $\Phi$, il existe un
entier $n>0$ tel que
$\mathscr{J}_0^n(\mathscr{O}_X/\mathscr{J})=0$
(\hyperref[I.9.3.4-fr]{9.3.4}), ce qui signifie que
$\mathscr{J}_0^n\subset\mathscr{J}$.
\end{proof}

\begin{env}[10.8.3]
\label{I.10.8.3-fr}
Soit maintenant $\mathscr{F}$ un $\mathscr{O}_X$-Module cohérent~; pour
tout $\mathscr{J}\in\Phi$,
$\mathscr{F}\otimes_{\mathscr{O}_X}(\mathscr{O}_X/\mathscr{J})$ est un
$\mathscr{O}_X$-Module cohérent
(\hyperref[I.9.1.1-fr]{9.1.1}), de support contenu dans $X'$, et que nous
identifierons le plus souvent à sa restriction à $X'$. Lorsque
$\mathscr{J}$ parcourt $\Phi$, ces faisceaux forment un
\emph{système projectif} de faisceaux de groupes abéliens.
\end{env}

\begin{definition}[10.8.4]
\label{I.10.8.4-fr}
Étant donnés une partie fermée $X'$ d'un préschéma localement noethérien
$X$ et un $\mathscr{O}_X$-Module cohérent $\mathscr{F}$, on appelle
\emph{complété de $\mathscr{F}$ le long de $X'$} et on désigne par
$\mathscr{F}_{/X'}$ ou par $\widehat{\mathscr{F}}$ (lorsqu'aucune confusion
n'est possible) la restriction à $X'$ du faisceau

